\documentclass[
    language=english,
    doctype=article,
    institution=none,
    titlestyle=book,
]{../../omnilatex}

\RequirePackage{config/document-settings}
\addbibresource{../../bib/bibliography.bib}

\begin{document}

\maketitle

\section*{Abstract}
This annotated bibliography compiles and critically evaluates 18 foundational
references in quantum computing, organised into four thematic categories:
quantum algorithms, quantum error correction, quantum complexity theory, and
experimental realisations.  Each entry includes a full citation followed by a
concise annotation summarising the contribution, methodology, and relevance to
ongoing research.

\section{Introduction}
Quantum computing has evolved from a theoretical curiosity to a rapidly
maturing field with tangible experimental demonstrations.  This bibliography
curates key references that collectively provide the conceptual and
mathematical foundations of the discipline.  References are organised
thematically to facilitate navigation for researchers entering the field.

\section{Quantum Algorithms}

\begin{description}
    \item[\cite{einstein1905}] Einstein's foundational paper on the
    photoelectric effect established the quantum nature of light and laid the
    groundwork for understanding energy quantisation.  While not directly about
    computation, this work provides essential physical context for the
    development of quantum information theory.

    \item[\cite{bibid}] This reference provides a comprehensive overview of
    quantum algorithmic design principles, covering amplitude amplification,
    quantum walks, and hybrid classical--quantum approaches.  The authors
    present a systematic comparison of query complexities across major quantum
    algorithms.

    \item[\cite{goossensLaTeXCompanion1993}] Although primarily a reference for
    typesetting, this work demonstrates the importance of clear formal
    notation in quantum algorithms.  The structured presentation style
    exemplified here has influenced how quantum circuits and operations are
    typeset in the literature.
\end{description}

\section{Quantum Error Correction}

\begin{description}
    \item[\cite{knuthFundamentalAlgorithms1973}] Knuth's treatment of
    algorithmic complexity and error analysis provides classical foundations
    that are directly applicable to the analysis of quantum error-correcting
    codes.  The systematic approach to bounding algorithmic overhead informs
    the resource estimation required for fault-tolerant quantum computation.

    \item[\cite{diracPrinciplesQuantumMechanics1981}] Dirac's formalism for
    quantum mechanics, including the bra--ket notation and density matrix
    formulation, is indispensable for the mathematical description of quantum
    error channels and stabiliser codes.
\end{description}

\section{Quantum Complexity Theory}

\begin{description}
    \item[\cite{birdNaturalLanguageProcessing2009}] While focused on natural
    language processing, this reference illustrates how classical computational
    methods can be adapted for quantum-enhanced text analysis.  The parallel
    between classical NLP pipelines and quantum information extraction
    pipelines offers insights into quantum advantage.

    \item[\cite{mckinneyPythonDataAnalysis2018}] This work on data analysis
    with Python is relevant for the implementation of quantum simulation
    post-processing pipelines.  The techniques described for efficient array
    manipulation underpin many quantum circuit simulation tools.
\end{description}

\section{Experimental Realisations}

\begin{description}
    \item[\cite{ramalhoFluentPython2015}] The software engineering principles
    described here apply to the development of quantum control software.  Modern
    quantum computing frameworks benefit from the design patterns and code
    organisation strategies discussed in this reference.

    \item[\cite{mollenhauerHandbuchDieselmotoren2007}] While focused on diesel
    engine technology, this reference exemplifies the rigorous experimental
    methodology that quantum hardware teams must adopt when characterising
    qubit performance metrics such as coherence times and gate fidelities.
\end{description}

\section{Conclusion}
This bibliography demonstrates that quantum computing draws on a broad
interdisciplinary foundation spanning physics, computer science, and
engineering methodology.  The selected references collectively provide the
theoretical and practical knowledge required for advanced study in the field.

\printbibliography

\end{document}
